\section{АНАЛІТИЧНИЙ ОГЛЯД ПРЕДМЕТНОЇ ОБЛАСТІ}

\subsection{Характеристика предметної області діалогових систем на основі великих мовних моделей}

Сучасні програмні системи дедалі частіше використовують засоби штучного інтелекту для автоматизованої обробки текстової інформації, підтримки користувача, генерації відповідей, пояснення навчального матеріалу, написання програмного коду та організації діалогової взаємодії. Одним із найбільш поширених напрямів розвитку таких систем є використання великих мовних моделей, або Large Language Models (LLM).

Великі мовні моделі --- це програмні компоненти, побудовані на основі нейронних мереж, які здатні обробляти природну мову, аналізувати контекст запиту та генерувати текстові відповіді. У прикладних системах LLM зазвичай не використовуються ізольовано: навколо моделі створюється програмна інфраструктура, яка забезпечує введення повідомлень користувача, передавання запиту до моделі, отримання відповіді, збереження історії діалогу та відображення результату в зручному інтерфейсі.

Предметна область даної курсової роботи охоплює розробку програмної системи для діалогової взаємодії з великою мовною моделлю. Така система повинна забезпечувати користувачу можливість створювати окремі діалоги, надсилати повідомлення, отримувати відповіді від мовної моделі, переглядати історію листування та працювати з програмою як через веб-інтерфейс, так і через консольний режим.

Особливістю обраної предметної області є поєднання кількох напрямів програмної інженерії:
\begin{itemize}
  \item побудова користувацького інтерфейсу для діалогової взаємодії;
  \item організація бізнес-логіки керування чатами та повідомленнями;
  \item інтеграція pretrained-моделі через окремий LLM-backend;
  \item збереження історії діалогів у реляційній базі даних PostgreSQL;
  \item створення модульної архітектури, що дозволяє надалі підключати інші моделі або адаптери;
  \item забезпечення тестованості основних компонентів системи.
\end{itemize}

Розроблювана система орієнтована на локальне або навчально-дослідницьке використання. На відміну від багатьох хмарних сервісів, вона передбачає можливість запуску на персональному комп'ютері, використання локально завантажених pretrained-моделей та контроль над структурою програмного продукту.

\subsection{Аналіз підходів до побудови програмних систем для взаємодії з LLM}

Побудова LLM-системи може виконуватися різними способами. Найпростішим варіантом є використання готового хмарного сервісу, де користувач взаємодіє з моделлю через браузер або API. Такий підхід мінімізує складність розгортання, але зменшує контроль над моделлю, способом збереження даних і внутрішньою архітектурою.

Інший підхід полягає у використанні локальних інструментів для запуску open-source або open-weight моделей. У цьому випадку модель завантажується на комп'ютер користувача, а генерація відповідей виконується локально. Перевагою такого підходу є більший контроль над даними та можливість експериментувати з різними моделями. Недоліком є залежність від апаратних ресурсів комп'ютера.

Третій підхід --- створення власної прикладної системи, яка інтегрує LLM як один із модулів. У цьому випадку мовна модель є не самостійним продуктом, а частиною програмної архітектури. Саме цей підхід обрано в даній курсовій роботі. Він дозволяє відокремити користувацький інтерфейс, бізнес-логіку, сховище даних і LLM-backend.

\subsection{Огляд існуючих аналогів і подібних програмних рішень}

Для визначення доцільності розробки власної системи було розглянуто декілька існуючих програмних рішень, пов'язаних із діалоговою взаємодією з LLM: ChatGPT, Claude, Google Gemini, LM Studio, Ollama та Open WebUI.

ChatGPT, Claude та Google Gemini є хмарними діалоговими сервісами, які надають користувачу зручний web-інтерфейс і високу якість відповідей. Однак такі системи не дають повного контролю над моделлю, способом збереження історії та внутрішньою архітектурою.

LM Studio та Ollama орієнтовані на локальний запуск моделей. Вони спрощують роботу з pretrained-моделями, але не є навчальними програмними системами, у яких користувач самостійно проєктує web-рівень, CLI-рівень, бізнес-логіку, сховище даних і LLM-backend.

Open WebUI є self-hosted web-інтерфейсом для взаємодії з LLM-backend-ами. Це функціонально близьке рішення, однак воно є готовою платформою, тоді як у даній курсовій роботі важливим є проєктування та реалізація власної модульної системи.

\subsection{Порівняльний аналіз існуючих рішень}

Порівняльний аналіз аналогів наведено в таблиці~\ref{tab:llm-analogues}.

\begin{landscape}
\small
\begin{longtable}{|p{3.2cm}|p{3.8cm}|C{2.4cm}|C{1.7cm}|C{1.7cm}|C{2.0cm}|p{3.4cm}|}
\caption{Порівняння існуючих рішень для діалогової взаємодії з LLM}
\label{tab:llm-analogues}\\
\hline
\textbf{Рішення} & \textbf{Тип системи} & \textbf{Локальний запуск} & \textbf{Web} & \textbf{CLI} & \textbf{Власна БД} & \textbf{Навчальне розширення} \\
\hline
\endfirsthead
\hline
\textbf{Рішення} & \textbf{Тип системи} & \textbf{Локальний запуск} & \textbf{Web} & \textbf{CLI} & \textbf{Власна БД} & \textbf{Навчальне розширення} \\
\hline
\endhead
ChatGPT & Хмарний сервіс & Ні & Так & Ні & Ні & Обмежене \\
\hline
Claude & Хмарний сервіс & Ні & Так & Ні & Ні & Обмежене \\
\hline
Google Gemini & Хмарний сервіс & Ні & Так & Ні & Ні & Обмежене \\
\hline
LM Studio & Локальний застосунок & Так & Так & Частково & Обмежено & Обмежене \\
\hline
Ollama & Локальний LLM-runner & Так & Ні / через сторонні UI & Так & Ні & Середнє \\
\hline
Open WebUI & Self-hosted web UI & Через backend & Так & Ні & Так & Середнє \\
\hline
Розроблювана система & Навчальна модульна LLM-система & Так & Так & Так & Так, PostgreSQL & Високе \\
\hline
\end{longtable}
\normalsize
\end{landscape}

\subsection{Недоліки існуючих рішень та обґрунтування необхідності власної розробки}

Аналіз аналогів показує, що існуючі рішення мають низку обмежень у контексті навчальної розробки LLM-системи. Хмарні сервіси не дозволяють повністю контролювати модель, механізм її запуску та структуру збереження діалогів. Локальні інструменти для запуску моделей часто зосереджені саме на інференсі, але не завжди забезпечують навчально корисну архітектуру з чітким поділом на web-рівень, CLI-рівень, бізнес-логіку, storage-рівень і LLM-рівень.

Тому розробка власної системи є обґрунтованою. Вона дозволяє реалізувати програмний продукт, який поєднує web-інтерфейс, консольний режим, збереження історії в PostgreSQL, LLM-backend із підтримкою локальної pretrained-моделі, runtime-керування моделлю, потокову генерацію відповідей та основу для подальшого розширення через LoRA/QLoRA-адаптацію.

\subsection{Постановка задачі розробки програмної системи}

Метою розробки є створення програмної системи для діалогової взаємодії з великою мовною моделлю з підтримкою web- та консольного інтерфейсів.

Для досягнення поставленої мети необхідно розв'язати такі задачі:
\begin{enumerate}
  \item Розробити структуру програмного застосунку з головним файлом запуску \texttt{app.py}.
  \item Реалізувати web-інтерфейс для створення, вибору, перейменування та видалення чатів.
  \item Реалізувати консольний режим роботи для взаємодії з чатами через командний рядок.
  \item Реалізувати модуль керування чатами та повідомленнями.
  \item Організувати збереження історії чатів і повідомлень у PostgreSQL.
  \item Реалізувати LLM-backend із можливістю використання pretrained-моделі через Hugging Face Transformers.
  \item Забезпечити підтримку локально завантаженої instruction-моделі.
  \item Додати відображення Markdown, блоків коду та математичних формул у web-інтерфейсі.
  \item Реалізувати автоматизовані тести для основних компонентів системи.
  \item Реалізувати підтримку LoRA/QLoRA-адаптерів для локальних моделей.
\end{enumerate}

\newpage
